% Generated draft: requires mathematical and source review.
\documentclass[11pt]{article}
\usepackage[margin=25mm]{geometry}
\usepackage{fontspec}
\setmainfont{Latin Modern Roman}
\usepackage{amsmath,amssymb,mathtools,unicode-math}
\setmathfont{Latin Modern Math}
\usepackage{enumitem,needspace,xurl,hyperref}
\setlength{\parindent}{0pt}
\setlength{\parskip}{5pt}
\setlength{\emergencystretch}{3em}
\newcommand{\N}{\mathbb{N}}
\newcommand{\Z}{\mathbb{Z}}
\newcommand{\Q}{\mathbb{Q}}
\newcommand{\R}{\mathbb{R}}
\newcommand{\C}{\mathbb{C}}
\newcommand{\F}{\mathbb{F}}
\newcommand{\A}{\mathbb{A}}
\newcommand{\PP}{\mathbb{P}}
\newcommand{\E}{\mathbb{E}}
\newcommand{\eps}{\varepsilon}
\newcommand{\dd}{\,\mathrm{d}}
\begin{document}
\section*{UnsolvedMath ID: 30000470}
\textbf{Source problem code:} catalogue entry
\subsection{Definitions and mathematical statement}
Are large cardinals consistent with stronger forms of Condensation?

This is a request to determine compatibility between specified large-cardinal axioms and specified strengthenings of a condensation principle. The supplied record does not identify either side precisely. It names Strong Condensation in the assessment, but gives no definition of that principle, its underlying hierarchy, its quantification over elementary substructures, or the ambient model in which it is to hold. It also supplies no particular large-cardinal notion, number of such cardinals, or comparison principle relative to which ``stronger'' is measured.

In the usual set-theoretic interpretation, one would formulate a consistency or relative-consistency statement over a fixed background theory, commonly ZFC, asserting a model in which the chosen large-cardinal assumptions and chosen condensation principle both hold, possibly assuming the consistency of stronger hypotheses. The source does not select ZFC explicitly, specify the required strength of the consistency assumption, or ask for an exact equiconsistency result. None of these is silently supplied.

Ordinary condensation for a particular inner-model hierarchy, a local condensation scheme, and a global principle called Strong Condensation are not interchangeable. Likewise, compatibility with one large-cardinal notion does not establish compatibility with every stronger notion. The supplied assessment reports both incompatibilities at some strengths and consistency for weaker principles with substantial large cardinals, but names neither the exact axioms nor the thresholds. Thus this record is a broad compatibility question whose precise instances and unresolved range must be recovered from the report; it is not a well-defined single universal consistency assertion.
\subsection{Short English statement}
The source asks how strong condensation principles can coexist with large cardinals. It specifies neither the condensation scheme nor the cardinal axioms or background consistency assumptions, so the exact compatibility problem remains unresolved at the formulation level.
\subsection{Sources}
\begin{itemize}
\item Local UnsolvedMath registry, permanent ID 30000470, code catalogue entry, ``Strong Condensation with Large Cardinals.'' Proposed year 2006; proposer unspecified.
\item \emph{Feinstrukturtheorie und Innere Modelle} (2007), Oberwolfach Reports, pp.\ 1215--1244, DOI 10.4171/owr/2006/20; \url{https://doi.org/10.4171/owr/2006/20}. The displayed report year 2007 differs from the DOI and proposed year 2006.
\item Supplied reference ``Friedman--Holy, Condensation and large cardinals'': \url{https://eudml.org/doc/286512}.
\item Statement status exact; the supplied verification says the statement was confirmed without textual change. This confirms wording, not the missing definitions.
\item Literature triage checked 2026-08-21: partially solved. It says Strong Condensation conflicts with some large-cardinal strengths, weaker condensation principles coexist with substantial large cardinals, and an unrestricted stronger-form consistency question remains open. No exact principles or strengths are given. Created 2026-08-25; updated 2026-09-16. These dated claims require independent verification.
\item Registry curation by Ulam AI / UnsolvedMath Contributors, CC BY 4.0; underlying report and literature provenance remain as supplied. Unreviewed statement-only draft.
\end{itemize}
\end{document}
